\begin{center}

\textbf{Partie A : études de deux fonctions}

\end{center}

On considère les deux fonctions $f$ et $g$ définies sur l'intervalle $[0~;~+\infty[$ par:

\[f(x) = 0,06\left(-x^2 +13,7x\right)\quad  \text{et}\quad  g(x) = (-0,15x + 2,2)\text{e}^{0,2x} - 2,2.\]

On admet que les fonctions $f$ et $g$ sont dérivables et on note $f'$ et $g'$ leurs fonctions dérivées respectives.

\medskip

\begin{enumerate}
\item On donne le tableau de variations complet de la fonction $f$ sur l'intervalle $[0~;~+\infty[$.

\begin{center}
\psset{unit=1cm,arrowsize=2pt 3}
\begin{pspicture}(6.5,2)
\psframe(6.5,2)\psline(0,1.5)(6.5,1.5)\psline(1.5,0)(1.5,2)
\uput[u](0.75,1.4){$x$} \uput[u](1.6,1.4){$0$} \uput[u](4,1.4){$6,85$} \uput[u](6,1.4){$+ \infty$} 
\rput(0.75,0.75){$f(x)$}\uput[u](1.65,0){$0$}\uput[d](4,1.5){$f(6,85)$}\uput[u](6,0){$- \infty$}
\psline{->}(1.75,0.25)(3.4,1.25)\psline{->}(4.6,1.25)(5.9,0.35)
\end{pspicture}
\end{center}

	\begin{enumerate}
		\item Justifier la limite de $f$ en $+\infty$.
		\item Justifier les variations de la fonction $f$.
		\item Résoudre l'équation $f(x) = 0$.
\end{enumerate}
\item 
	\begin{enumerate}
		\item Déterminer la limite de $g$ en $+\infty$.
		\item Démontrer que, pour tout réel $x$ appartenant à $[0~;~+\infty[$ on a : $g'(x) = (- 0,03x + 0,29)\text{e}^{0,2x}$.
		\item Étudier les variations de la fonction $g$ et dresser son tableau de variations sur $[0~;~+\infty[$.
		
Préciser une valeur approchée à $10^{-2}$ près du maximum de $g$.
		\item Montrer que l'équation $g(x) = 0$ admet une unique solution non nulle et déterminer, à $10^{-2}$ près, une valeur approchée de cette solution.
	\end{enumerate}
\end{enumerate}

\bigskip

\textbf{Partie B : trajectoires d'une balle de golf}

\medskip

Pour frapper la balle, un joueur de golf utilise un instrument appelé \og club\fg{} de golf.

On souhaite exploiter les fonctions $f$ et $g$ étudiées en Partie A pour modéliser de deux façons différentes la trajectoire d'une balle de golf. On suppose que le terrain est parfaitement plat.

On admettra ici que $13,7$ est la valeur qui annule la fonction $f$ et une approximation de la valeur qui annule la fonction $g$.

On donne ci-dessous les représentations graphiques de $f$ et $g$ sur l'intervalle [0~;~13,7].

\begin{center}
\psset{unit=1cm}
\begin{pspicture*}(-1,-1)(14,3.5)
\psgrid[gridlabels=0pt,subgriddiv=1,gridwidth=0.2pt]
\psaxes[linewidth=1.25pt,Dx=20,Dy=20]{->}(0,0)(-1,-1)(14,3.5)
\psplot[plotpoints=1000,linewidth=1.25pt,linecolor=red]{0}{13.7}{13.7 x mul x dup mul sub 0.06 mul}\uput[ul](3,2){\red $\mathcal{C}_f$}
\psplot[plotpoints=1000,linewidth=1.25pt,linecolor=blue]{0}{13.7}{2.2 0.15 x mul sub 2.71828 0.2 x mul exp mul 2.2 sub}\uput[ur](12,2.2){\blue $\mathcal{C}_g$}
\uput[dl](0,0){0}\uput[d](1,0){1}\uput[l](0,1){1}\uput[d](13.7,0){13,7}
\end{pspicture*}
\end{center}

Pour $x$ représentant la distance horizontale parcourue par la balle en dizaine de yards après la frappe, (avec $0 < x < 13,7$), $f(x)$ (ou $g(x)$ selon le modèle) représente la hauteur correspondante de la balle par rapport au sol, en dizaine de yards (1 yard correspond à environ $0,914$ mètre).

On appelle \og angle de décollage \fg{} de la balle, l'angle entre l'axe des abscisses et la tangente à la courbe ($\mathcal{C}_f$ ou $\mathcal{C}_g$ selon le modèle) en son point d'abscisse $0$. Une mesure de l'angle de décollage de la balle est
un nombre réel $d$ tel que $\tan (d)$ est égal au coefficient directeur de cette tangente.

De même, on appelle \og angle d'atterrissage \fg{} de la balle, l'angle entre l'axe des abscisses et la tangente
à la courbe ($\mathcal{C}_f$ ou $\mathcal{C}_g$ selon le modèle) en son point d'abscisse $13,7$. Une mesure de l'angle d'atterrissage de la balle est un nombre réel $a$ tel que $\tan (a)$ est égal à l'opposé du coefficient directeur de cette tangente.

Tous les angles sont mesurés en degré.

\begin{center}
\begin{tabularx}{\linewidth}{|p{5.75cm}|X|}\hline
Le schéma illustre les angles de décollage et d'atterrissage associés à la courbe 
$\mathcal{C}_f$&Le schéma illustre les angles de décollage et d'atterrissage
associés à la courbe $\mathcal{C}_g$.\\ \hline
\psset{unit=0.4cm}
\begin{pspicture}(-0.3,-0.75)(14,3.2)

\psaxes[linewidth=1.25pt,Dx=20,Dy=10](0,0)(-0.3,-0.3)(14,3)
\psplot[plotpoints=1000]{0}{2}{0.822 x mul}
\psplot[plotpoints=1000,linewidth=1.25pt,linecolor=red]{0}{13.7}{13.7 x mul x dup mul sub 0.06 mul}\uput[ul](3,2){\red $\mathcal{C}_f$}
\psline(13.7,0)(11,2)
\psarc(0,0){0.7}{0}{40}\rput(1.2,0.5){\footnotesize $d$}
\psarc(13.7,0){0.7}{140}{180}\rput(12.4,0.5){\footnotesize $a$}
\uput[d](13.7,0){\footnotesize 13,7}
\end{pspicture}&\psset{unit=0.5cm}
\begin{pspicture}(-0.3,-0.75)(14,3.2)

\psaxes[linewidth=1.25pt,Dx=20,Dy=10](0,0)(-0.3,-0.5)(14,3)
\psplot[plotpoints=1000,linewidth=1.25pt,linecolor=blue]{0}{13.7}{2.2 0.15 x mul sub 2.71828 0.2 x mul exp mul 2.2 sub}
\psline(0,0)(10.4,2.8)
\psline(13.7,0)(12,2.9)
\psarc(0,0){1}{0}{20}\rput(2,0.25){\footnotesize $d$}
\psarc(13.7,0){1}{120}{180}\rput(12.4,0.5){\footnotesize $a$}
\uput[d](13.7,0){\footnotesize 13,7}
\end{pspicture}\\ \hline
\end{tabularx}
\end{center}

\medskip

\begin{enumerate}
\item \emph{Première modélisation} : on rappelle qu'ici, l'unité étant la dizaine de yards, $x$ représente la distance horizontale parcourue par la balle après la frappe et $f(x)$ la hauteur correspondante de la balle.

Selon ce modèle :
	\begin{enumerate}
		\item Quelle est la hauteur maximale, en yard, atteinte par la balle au cours de sa trajectoire ?
		\item Vérifier que $f'(0) = 0,822$.
		\item Donner une mesure en degré de l'angle de décollage de la balle, arrondie au dixième. (On pourra éventuellement utiliser le tableau ci-dessous).
		\item Quelle propriété graphique de la courbe $\mathcal{C}_f$ permet de justifier que les angles de décollage et d'atterrissage de la balle sont égaux ?
	\end{enumerate}
\item \emph{Seconde modélisation } : on rappelle qu'ici, l'unité étant la dizaine de yards, $x$ représente la distance horizontale parcourue par la balle après la frappe et $g(x)$ la hauteur correspondante de la balle.

Selon ce modèle :

	\begin{enumerate}
		\item Quelle est la hauteur maximale, en yard, atteinte par la balle au cours de sa trajectoire ?

On précise que $g'(0) = 0,29$ et $g'(13,7) \approx -1,87$.
		\item Donner une mesure en degré de l'angle de décollage de la balle, arrondie au dixième. (On pourra éventuellement utiliser le tableau ci-dessous).
		\item Justifier que $62$ est une valeur approchée, arrondie à l'unité près, d'une mesure en degré de l'angle d'atterrissage de la balle.
	\end{enumerate}
\medskip

\textbf{Tableau :} extrait d'une feuille de calcul donnant une mesure en degré d'un angle quand on connait sa tangente :

\begin{center}
\begin{tabularx}{\linewidth}{|c|*{13}{>{\centering \arraybackslash}X|}}\hline
&A	&B	&C	&D	&E	&F	&G	&H	&I	&J	&K	&L	&M\\ \hline
1&$\tan (\theta)$&0,815&0,816&0,817&0,818&0,819&0,82&0,821&0,822&0,823&0,824&0,825&0,826\\ \hline
2&\scriptsize $\theta$ en degrés&39,18&39,21&39,25&39,28&39,32&39,35&39,39&39,42&39,45&39,49&39,52&39,56\\ \hline
3&	&&&&&&&&&&&&\\ \hline
4&$\tan (\theta)$&0,285& 0,286& 0,287& 0,288& 0,289&0,29&0,291&0,292&
0,293&0,294&0,295 &0,296\\ \hline
5&\scriptsize  $\theta$ en degrés&15,91 &15,96&16,01& 16,07& 16,12& 16,17& 16,23& 16,28& 16,33& 16,38& 16,44& 16,49\\ \hline
\end{tabularx}
\end{center}
\end{enumerate}

\bigskip

\begin{center}
\textbf{Partie C : interrogation des modèles}
\end{center}

À partir d'un grand nombre d'observations des performances de joueurs professionnels, on a obtenu les résultats moyens suivants:

\begin{center}
\begin{tabularx}{\linewidth}{|*{4}{>{\centering \arraybackslash \footnotesize}X|}}\hline
Angle de décollage en degré&Hauteur maximale en yard&Angle d'atterrissage en degré&Distance horizontale en yard au point de chute\\ \hline
24&32&52&137\\ \hline
\end{tabularx}
\end{center}

Quel modèle, parmi les deux étudiés précédemment, semble le plus adapté pour décrire la frappe de la balle par un joueur professionnel ? La réponse sera justifiée.
